\scanpage{13}

On obtient donc des bijections inverses l'une de l'autre.
\diagram{\begin{tikzcd}
\pi_0(U\otimes\Lambda) \arrow[r,bend left=15,"\exp"] & \gamma_0(U\otimes_k\Lambda) \arrow[l,bend left=15,"\log"]
\end{tikzcd}}
où les indices $0$ indiquent qu'on
prend les primitifs resp\textsuperscript{t} multiplicatifs
``relatifs'' à l'augmentation $U\otimes\Lambda\longrightarrow U$.

[Par passage à la limite projective \inserted{en $\Lambda$}, on déduit
un résultat analogue si $\Lambda$ est une
limite projective de $k$-algèbres infinitésimales].

Notons que \deleted{si $\Lambda$ est plat, \unsure{--- si}}
l'on a d'ailleurs un
hom canonique de $\Lambda$-modules
\diagram{\begin{tikzcd}
\pi(U)\otimes_k I(\Lambda) \arrow[r,"\alpha"] & \pi_0(U\otimes_k\Lambda)
\end{tikzcd}}
qui est un isom si $\Lambda$ est $k$-plat,
ou si
\[
0\;\text{---}\;\deleted{\unsure{\mathfrak g}}\pi(U)
\;\text{---}\;\unsure{I(U)}\xrightarrow{\delta}I(U)\otimes I(U)
\]
est une suite \uuline{universellement exacte}. On
en déduit donc un hom canonique
\diagram{\begin{tikzcd}
\pi(U)\otimes_k I(\Lambda) \arrow[r,"\exp\circ\alpha"] & 1+\pi_0(U\otimes_k\Lambda)
\end{tikzcd}}
qui est un isom. dans les cas
\unsure{qu'on vient de préciser}.

\margin{[i.e. $\Delta$ est comm. et assoc,
\deleted{$\unsure{A\longrightarrow A}$}
$I(U)=\bigcup U_i$, avec $\delta U_i\subset U_i\otimes U_i$ et
$U_i$ de dim. finie, enfin $\gamma(U)=\{e\}$]}

\margin{La donnée d'une \unsure{telle} est connue quand on connaît
$\mathfrak g=\pi(U)$\deleted{$\unsure{\otimes}$}
(formule de Hausdorff-Campbell), donc
(comme par \unsure{Poincaré--B.W.} $\mathfrak g\simeq\pi(U(\mathfrak g)))$
on voit que $U\simeq U(\mathfrak g)$ avec $\mathfrak g\simeq\pi(U)$ :
\uline{th. de Cartier} $\Longleftarrow$ \unsure{précédent}.
Si $\mathfrak g$ est de dim. finie, le dual de $\unsure{U}$ est un
\uline{anneau local} \uline{régulier}. \uline{Cor} Un groupe alg. en car $0$
est \unsure{toujours lisse}.}

Conséquence : Supposons $k$ corps,
\inserted{$U$ une hyperalgèbre \uline{de Cartier}}.
Alors $U$ (avec $\Delta$ \uline{comm.} et \uline{ass.}) est connu
quand on connaît le foncteur $\pi_0(U\otimes_k\Lambda)$.

\mspage{12}
